\section{Structural properties and witnessed perfect play}
\label{sec:properties}

A property query must specify the mathematical predicate before selecting an
analysis method. Structural predicates inspect rational verifier data: acceptance
rewards, Boolean winning conditions, and dependence only on the parity of binary
answers. An XOR score can have constant or signed parity rewards; an XOR winning
predicate accepts exactly one parity at each selected question. These notions
apply at arbitrary party count and do not themselves license the two-party vector
SDP. The question domain is explicit: all stored inputs or only the positive-mass
sampling support.

Perfect play concerns the sampled game and requires acceptance rewards on that
support. It is an attainment statement, not merely the equality of a supremum
with one. The initial property fragment considers ordinary local and finite
tensor-product quantum strategies. More restrictive resources and self-testing
claims require additional rules rather than an inferred Boolean result.

\begin{proposition}[Sufficient property evidence]
For an acceptance game, a verified finite quantum strategy with score one proves
finite quantum perfectness, while a checked quantum upper bound below one
excludes it. A verified perfect local strategy also gives a perfect finite quantum
strategy by using one-dimensional systems and identity/zero effects. A perfect
finite quantum strategy together with a checked local upper bound below one proves
pseudotelepathy. Either a perfect local strategy or exclusion of quantum perfectness
refutes pseudotelepathy.
\end{proposition}
\begin{proof}
The point certificates establish admissible attained scores directly. An upper
bound below one excludes every score-one realization in its stated model.
Deterministic local responses have the indicated finite quantum embedding.
Pseudotelepathy is the conjunction of finite quantum perfectness and the absence
of a perfect local strategy, so the remaining conclusions follow.
\end{proof}
A nonsignalling score-one point is not a quantum point. A generic lower bound of
one is likewise not an attainment certificate; nonclosed quantum correlation sets
make the distinction substantive \cite{slofstra2019}. Numerical rank and optimizer
uniqueness do not repair this missing implication. The property engine therefore
records checked attainment as a separate fact and keeps its originating resource.

Every admitted attained score also supplies a lower value bound. This fact is
included in consistency checking even if a future producer omits a separate lower
bound field. For acceptance games, checked lower bounds above one, upper bounds
below zero, or attained scores outside $[0,1]$ are contradictory premises and
quarantine the assessment. Loose bounds outside that interval in the opposite
direction are not contradictions. Absence of a proof remains unknown; it is not
evidence that the property is false.

All property definitions bind to the exact verifier and support convention.
Structural counterexamples record locations and reasons that can be replayed on
the input data. Evidence-backed assessments recheck raw artifacts through the same
admission layer as value claims. Invalid, unsupported, exhausted, or conflicting
checks produce operational outcomes separate from mathematical verdicts. Receipts
are not standalone proof tokens, and no general strategy-discovery procedure is
launched implicitly by property assessment.
